\setcounter{page}{76}
\textbf{Definition.}
Let \(\cat{A}\) be an abelian category, and let \(X^\bullet\) be a complex over \(\cat{A}\).
We refer to [M, Ch.\ IV, \S4] for the definition of double complexes, morphisms of double complexes, and homotopies between such morphisms.

A \textit{Cartan--Eilenberg resolution} of \(X^\bullet\) is a double complex \(C^{\cdot\cdot}\) of injective objects of \(\cat{A}\) with \(C^{p,q}=0\) for \(q<0\), together with an augmentation
\(\epsilon\colon X^\bullet\to C^{\cdot,0}\), such that for every \(p\in\mathbb{Z}\), the induced maps
\[
B^p(\epsilon)\colon B^p(X^\bullet)\to B_\mathrm{I}^p(C^{\cdot\cdot}),\qquad
H^p(\epsilon)\colon H^p(X^\bullet)\to H_\mathrm{I}^p(C^{\cdot\cdot})
\]
are injective resolutions in the ordinary sense (cf.\ [M, Ch.\ XVII, \S], where this is called an injective resolution of the complex \(X^\bullet\)).

We recall standard properties of double complexes and Cartan--Eilenberg resolutions.

\textbf{Lemma 7.5.}
\begin{enumerate}
\item[(a)] If \(\cat{A}\) has enough injectives, every complex \(X^\bullet\) admits a Cartan--Eilenberg resolution.
\item[(b)] Let \(f\colon X^\bullet\to Y^\bullet\) be a morphism of complexes. If \(C^{\cdot\cdot},D^{\cdot\cdot}\) are Cartan--Eilenberg resolutions of \(X^\bullet,Y^\bullet\) respectively, there exists a double complex morphism \(F\colon C^{\cdot\cdot}\to D^{\cdot\cdot}\) lifting \(f\).
\item[(c)] If \(f,g\colon X^\bullet\to Y^\bullet\) are homotopic, and \(F,G\) lift \(f,g\) respectively to Cartan--Eilenberg resolutions, then \(F\) is homotopic to \(G\).
\end{enumerate}

\setcounter{page}{77}
\begin{enumerate}
\item[(d)] Homotopic double complex morphisms induce homotopic morphisms of their associated simple complexes.
\item[(e)] Let \(\sigma_{\le n}^\mathrm{II}\) denote truncation with respect to the second index. If \(F,G\) are homotopic double complex maps, their restrictions to \(\sigma_{\le n}^\mathrm{II}(C^{\cdot\cdot})\to\sigma_{\le n}^\mathrm{II}(D^{\cdot\cdot})\) are also homotopic.
\item[(f)] Let \(\epsilon\colon X^\bullet\to C^{\cdot\cdot}\) be an augmentation with \(C^{p,q}=0\) for \(q<0\) and \(q>n_0\), such that each \(X^p\to C^{p,\bullet}\) is a resolution. Then the natural map \(X^\bullet\to s(C^{\cdot\cdot})\) to the associated simple complex is a quasi-isomorphism.
\end{enumerate}

\textit{Proofs.}
(a), (b), (c): [M, Ch.\ XVII, Prop.\ 1.2].
(d): [M, Ch.\ IV, \S4].
(e): Elementary.
(f): Follows from the spectral sequence of a double complex (EGA O\(_{\mathrm{III}}\) 11.3.3(ii)) or direct verification.

\textit{Proof of Proposition 7.4 (continued).}
Given \(X^\bullet\in L\), choose a Cartan--Eilenberg resolution \(C^{\cdot\cdot}\) (exists since \(\cat{A}\) has enough injectives). Define the truncated double complex \(C'^{\cdot\cdot}=\sigma_{\le n}^\mathrm{II}(C^{\cdot\cdot})\) by
\[
C'^{p,q}=
\begin{cases}
C^{p,q} & q<n,\\
\ker d_2^{p,n} & q=n,\\
0 & q>n.
\end{cases}
\]

\setcounter{page}{78}
For each \(p\in\mathbb{Z}\), we have an exact sequence
\[
0\to X^p\to C'^{p,0}\to C'^{p,1}\to\cdots\to C'^{p,n-1}\to C'^{p,n}\to 0.
\]
The terms \(C'^{p,0},\dots,C'^{p,n-1}\) are injective; \(X^p\in P\) implies \(C'^{p,n}\) is \(F\)-acyclic.
This resolution computes \(\mathbf{R}F(X^p)\), giving an isomorphism
\[
\varphi(t^p)\colon G(X^p)=R^nF(X^p) \to F(C'^{p,n})\big/\operatorname{im}F(C'^{p,n-1}),
\]
where \(t^p\colon X^p\to C^{p,\bullet}\) is the resolution map.
This induces a morphism \(\alpha^p\colon F(C^{p,n})\to G(X^p)\), and hence a double complex morphism
\[
\alpha\colon F(C'^{\cdot\cdot})\to G(X^\bullet),
\]
with \(G(X^\bullet)\) concentrated in the \(n\)-th row.

Passing to associated simple complexes:
\[
s(\alpha)\colon F\big(s(C'^{\cdot\cdot})\big)\to G(X^\bullet)[-n].
\]

\setcounter{page}{79}
By Lemma 7.5(f), the augmentation \(u\colon X^\bullet\to s(C^{\cdot\cdot})\) is a quasi-isomorphism into an \(F\)-acyclic complex, yielding an isomorphism
\[
\varphi(u)\colon \mathbf{R}F(X^\bullet)\to F\big(s(C^{\cdot\cdot})\big).
\]
Since \(X^\bullet\) consists of \(G\)-acyclic objects, we have
\[
\varphi(\id_{X^\bullet})\colon G(X^\bullet)\to \mathbf{L}G(X^\bullet).
\]
Composing these gives the desired morphism
\[
\psi(X^\bullet)\colon \mathbf{R}F(X^\bullet)\to \mathbf{L}G(X^\bullet).
\]

The morphism \(\psi(X^\bullet)\) is independent of the Cartan--Eilenberg resolution choice.
Using Lemma 7.5, one verifies compatibility with morphisms \(f\colon X^\bullet\to Y^\bullet\) in \(L\), so \(\psi\) is a functor morphism
\[
\psi\colon \mathbf{R}F\to \mathbf{L}G[-n].
\]

To see \(\psi\) is an isomorphism: \(\mathbf{R}F,\mathbf{L}G\) are way-out in both directions. By the way-out functor lemma, reduce to checking \(\psi(X)\) is an isomorphism for \(X\in P\), which holds by construction. q.e.d.

\setcounter{page}{80}
\textbf{Proposition 7.6.}
Let \(\cat{A}\) have enough injectives, \(X^\bullet\in K^+(\cat{A})\). The following are equivalent:
\begin{enumerate}
\item[(i)] \(X^\bullet\) admits a quasi-isomorphism into a bounded complex of injective objects of \(\cat{A}\).
\item[(ii)] The functor \(\mathbf{R}\mathcal{H}\textit{om}(\,\cdot\,,X^\bullet)\colon D(\cat{A})^\circ\to D(\mathbf{Ab})\) is way-out left (hence way-out both ways).
\item[(iii)] There exists \(n_0\in\mathbb{Z}\) such that \(\operatorname{Ext}^i(Y,X^\bullet)=0\) for all \(Y\in\operatorname{Ob}\cat{A}\) and all \(i>n_0\).
\end{enumerate}

\textit{Proof.}
\((i)\Rightarrow(ii)\): Assume \(X^\bullet\) is bounded injective. Then \(\mathbf{R}\mathcal{H}\textit{om}(Y^\bullet,X^\bullet)\cong\mathcal{H}\textit{om}(Y^\bullet,X^\bullet)\), which is clearly way-out left.

\((ii)\Rightarrow(iii)\): Since \(F=\mathbf{R}\mathcal{H}\textit{om}(\,\cdot\,,X^\bullet)\) is way-out left, choose \(n_0\) such that \(H^i(F(Y^\bullet))=0\) whenever \(H^i(Y^\bullet)=0\) for \(i<n_0\).
For fixed \(Y\in\cat{A}\), set \(Y'^\bullet = T^{-n_0}(Y)\). Then \(H^i(Y'^\bullet)=0\) for \(i<n_0\), so \(H^i(F(Y'^\bullet))=0\) for \(i>0\). But
\[
H^i(F(Y'^\bullet))=H^i\big(F(T^{-n_0}Y)\big)=H^{i+n_0}(F(Y)),
\]
hence \(\operatorname{Ext}^{i+n_0}(Y,X^\bullet)=0\) for all \(i>0\), i.e., \(\operatorname{Ext}^j(Y,X^\bullet)=0\) for \(j>n_0\).

\setcounter{page}{81}
\((iii)\Rightarrow(i)\): Take an injective resolution \(X'^\bullet\in K^+(\cat{A})\) of \(X^\bullet\) (Lemma 4.6). Suppose \(H^m(X'^\bullet)\neq0\) for some \(m>n_0\). From the exact sequence
\[
0\to B^m(X'^\bullet)\to Z^m(X'^\bullet)\to H^m(X'^\bullet)\to 0,
\]
the inclusion is proper. Choose \(Y=Z^m(X'^\bullet)\); then
\[
\Hom(Y,B^m(X'^\bullet))\subsetneq \Hom(Y,Z^m(X'^\bullet)).
\]
But \(\operatorname{Ext}^m(Y,X^\bullet)=0\) by hypothesis, which forces the map to be an isomorphism, a contradiction. Thus \(H^i(X'^\bullet)=0\) for all \(i>n_0\).

\setcounter{page}{82}
For \(n\ge n_0\), the truncation map \(i\colon\sigma_{\le n}(X'^\bullet)\to X'^\bullet\) is a quasi-isomorphism.
For \(n\ge n_0\), \(Z^{n+1}(X'^\bullet)=B^{n+1}(X'^\bullet)\) is injective.
Consider the exact sequence
\[
0\to\sigma_{\le n}(X'^\bullet)\to\tau_{\le n}(X'^\bullet)\to B^{n+1}(X'^\bullet)[-n]\to 0.
\]
We show \(\operatorname{Ext}^1(Y,B^{n+1})=0\) for all \(Y\in\cat{A}\). We have
\[
\operatorname{Ext}^1(Y,B^{n+1})
=\operatorname{Ext}^{n+1}\big(Y,B^{n+1}[-n]\big).
\]

\setcounter{page}{83}
From the Ext long exact sequence:
\[
\cdots\to\operatorname{Ext}^{n+1}\big(Y,\tau_{\le n}(X'^\bullet)\big)
\to\operatorname{Ext}^{n+1}\big(Y,B^{n+1}[-n]\big)
\to\operatorname{Ext}^{n+2}\big(Y,\sigma_{\le n}(X'^\bullet)\big)\to\cdots.
\]
The left term vanishes because \(\tau_{\le n}(X'^\bullet)\) is bounded injective; the right term vanishes by hypothesis. Hence \(\operatorname{Ext}^1(Y,B^{n+1})=0\), so \(B^{n+1}\) is injective.
Thus \(\sigma_{\le n}(X'^\bullet)\) is a bounded injective complex quasi-isomorphic to \(X^\bullet\).

\textbf{Definition.}
If \(X^\bullet\in K^+(\cat{A})\) satisfies the equivalent conditions of Proposition 7.6, we say \(X^\bullet\) has \textit{finite injective dimension}.

\setcounter{page}{84}
\textbf{Corollary 7.7.}
The complexes of finite injective dimension form a triangulated localizing subcategory of \(K(\cat{A})\).

\textit{Proof.}
Closure under triangles follows from the Ext long exact sequence. Localizing follows from Proposition 3.3: any complex quasi-isomorphic to a finite injective dimension complex also has finite injective dimension.

\textbf{Definition.}
Denote by \(K^+(\cat{A})_{\mathrm{fid}}\) the full subcategory of \(K^+(\cat{A})\) consisting of complexes of finite injective dimension. The corresponding derived subcategory is \(D^+(\cat{A})_{\mathrm{fid}}\).

\textbf{Remark.}
These subcategories differ from \(K_{\cat{A}'}(\cat{A}),D_{\cat{A}'}(\cat{A})\), as they cannot be characterized solely by conditions on the cohomology objects \(H^i(X^\bullet)\).